% GENERATED FILE. Edit canonical lessons, then run npm run generate:books.
% canonical-source-hash: fnv1a64:9cf37d7f7b7592f5
% canonical-lessons: BN-C438-monobhab, BN-C438-onubhuti, BN-C438-abeg, BN-C438-onubhob-kora, BN-C438-ter-paoya

\chapter{Mind and Feeling}
\label{ch:bn-mind-and-feeling}
\hlchaptermodality{438}

\begin{chapteropening}
\textbf{By the end of this chapter:} \emph{I can talk about the mind and about what I feel.}

Complete the last lesson of chapter 438: I can talk about the mind and about what I feel.
\end{chapteropening}

\section[\texorpdfstring{\bn{মনোভাব} (mônobhāb)}{মনোভাব (mônobhāb)}]{\bn{মনোভাব} (mônobhāb) — an attitude; a frame of mind}

\label{lesson:BN-C438-monobhab}



\begin{quote}
Before the new one: say the Bengali for a female singer, then the Bengali for a female dancer.
\end{quote}

\subsection*{You'll want to know: \bn{মনোভাব}}
\textbf{\bn{মনোভাব}} — \emph{mônobhāb} — \textquotedblleft{}an attitude; a frame of mind\textquotedblright{}.

\begin{quote}
\bn{নতুন} \bn{কাজ} \bn{নিয়ে} \bn{অরুণের} \bn{মনোভাব} \bn{ভালো।}
\end{quote}

\begin{quote}
\emph{nôtun kāj niye ôruner mônobhāb bhālo.}
\end{quote}

\begin{quote}
Arun has a good attitude to the new work.
\end{quote}

\begin{cousinweb}
\textbf{\bn{মনোভাব}} (\emph{mônobhāb}) joins Sanskrit \emph{manas}, \textquotedblleft{}mind\textquotedblright{} (CDIAL 9822), in its compound form \emph{mano-}, and \emph{bhāva}, \textquotedblleft{}state, disposition\textquotedblright{} (CDIAL 9475): the state of the mind. \emph{Manas} is the \bn{মনে} of \bn{মনে} \bn{রাখা} (taught), to remember.
\end{cousinweb}

\begin{grammarlens}[title={What you've built}]
One more word for describing what you lived and felt.
\end{grammarlens}

\subsection*{Your turn}
\begin{itemize}
\raggedright
  \item \emph{Say it:} \emph{mônobhāb}
  \item \emph{Say it:} \emph{mônobhāb}, once more
  \item \emph{Say it:} \emph{nôtun kāj niye ôruner mônobhāb bhālo.} — Arun has a good attitude to the new work.
  \item \emph{Recall:} say the Bengali for a female singer, then the Bengali for a female dancer, then say \emph{mônobhāb} again
\end{itemize}

\begin{tcolorbox}[breakable,colback=teal!4,colframe=teal!35!black,title={Before you move on}]
What does \bn{মনোভাব} mean? (\textquotedblleft{}An attitude; a frame of mind\textquotedblright{}.) Say it once more.
\end{tcolorbox}
\section[\texorpdfstring{\bn{অনুভূতি} (ônubhūti)}{অনুভূতি (ônubhūti)}]{\bn{অনুভূতি} (ônubhūti) — a feeling; a sensation}

\label{lesson:BN-C438-onubhuti}



\begin{quote}
Before the new one: say the Bengali for a female dancer, then the Bengali for an attitude.
\end{quote}

\subsection*{You'll want to know: \bn{অনুভূতি}}
\textbf{\bn{অনুভূতি}} — \emph{ônubhūti} — \textquotedblleft{}a feeling; a sensation\textquotedblright{}.

\begin{quote}
\bn{এই} \bn{গানটা} \bn{শুনে} \bn{আমার} \bn{একটা} \bn{নতুন} \bn{অনুভূতি} \bn{হল।}
\end{quote}

\begin{quote}
\emph{ei gānṭā shune āmār ekṭā nôtun ônubhūti hôlô.}
\end{quote}

\begin{quote}
Listening to this song, I had a new feeling.
\end{quote}

\begin{cousinweb}
\textbf{\bn{অনুভূতি}} (\emph{ônubhūti}) is a tatsama, Sanskrit \emph{anubhūti}, \textquotedblleft{}perception\textquotedblright{} (Monier-Williams); in Bengali it names a feeling or a sensation.
\end{cousinweb}

\begin{grammarlens}[title={What you've built}]
One more word for describing what you lived and felt.
\end{grammarlens}

\subsection*{Your turn}
\begin{itemize}
\raggedright
  \item \emph{Say it:} \emph{ônubhūti}
  \item \emph{Say it:} \emph{ônubhūti}, once more
  \item \emph{Say it:} \emph{ei gānṭā shune āmār ekṭā nôtun ônubhūti hôlô.} — Listening to this song, I had a new feeling.
  \item \emph{Recall:} say the Bengali for a female dancer, then the Bengali for an attitude, then say \emph{ônubhūti} again
\end{itemize}

\begin{tcolorbox}[breakable,colback=teal!4,colframe=teal!35!black,title={Before you move on}]
What does \bn{অনুভূতি} mean? (\textquotedblleft{}A feeling; a sensation\textquotedblright{}.) Say it once more.
\end{tcolorbox}
\section[\texorpdfstring{\bn{আবেগ} (ābeg)}{আবেগ (ābeg)}]{\bn{আবেগ} (ābeg) — emotion; strong feeling}

\label{lesson:BN-C438-abeg}



\begin{quote}
Before the new one: say the Bengali for an attitude, then the Bengali for a feeling.
\end{quote}

\subsection*{You'll want to know: \bn{আবেগ}}
\textbf{\bn{আবেগ}} — \emph{ābeg} — \textquotedblleft{}emotion; strong feeling\textquotedblright{}.

\begin{quote}
\bn{বাবা} \bn{আবেগে} \bn{কেঁদে} \bn{ফেললেন।}
\end{quote}

\begin{quote}
\emph{bābā ābege kẽde phellen.}
\end{quote}

\begin{quote}
Father was so moved that he burst into tears.
\end{quote}

\bn{আবেগে}, overcome with feeling.

\begin{cousinweb}
\textbf{\bn{আবেগ}} (\emph{ābeg}) is a tatsama, Sanskrit \emph{āvega}, \textquotedblleft{}haste born of excitement, agitation\textquotedblright{} (Monier-Williams); Bengali uses it for a strong feeling that carries you along.
\end{cousinweb}

\begin{grammarlens}[title={What you've built}]
One more word for describing what you lived and felt.
\end{grammarlens}

\subsection*{Your turn}
\begin{itemize}
\raggedright
  \item \emph{Say it:} \emph{ābeg}
  \item \emph{Say it:} \emph{ābeg}, once more
  \item \emph{Say it:} \emph{bābā ābege kẽde phellen.} — Father was so moved that he burst into tears.
  \item \emph{Recall:} say the Bengali for an attitude, then the Bengali for a feeling, then say \emph{ābeg} again
\end{itemize}

\begin{tcolorbox}[breakable,colback=teal!4,colframe=teal!35!black,title={Before you move on}]
What does \bn{আবেগ} mean? (\textquotedblleft{}Emotion; strong feeling\textquotedblright{}.) Say it once more.
\end{tcolorbox}
\section[\texorpdfstring{\bn{অনুভব} \bn{করা} (ônubhôb kôrā)}{অনুভব করা (ônubhôb kôrā)}]{\bn{অনুভব} \bn{করা} (ônubhôb kôrā) — to feel; to sense}

\label{lesson:BN-C438-onubhob-kora}



\begin{quote}
Before the new one: say the Bengali for a feeling, then the Bengali for emotion.
\end{quote}

\subsection*{You'll want to know: \bn{অনুভব} \bn{করা}}
\textbf{\bn{অনুভব} \bn{করা}} — \emph{ônubhôb kôrā} — \textquotedblleft{}to feel; to sense\textquotedblright{}.

\begin{quote}
\bn{মীরা} \bn{বন্ধুর} \bn{ব্যথা} \bn{অনুভব} \bn{করল।}
\end{quote}

\begin{quote}
\emph{mīrā bôndhur bæthā ônubhôb kôrlô.}
\end{quote}

\begin{quote}
Mira felt her friend's pain.
\end{quote}

Past (he/she …): \bn{তিনি} \bn{ভয়} \bn{অনুভব} \bn{করলেন}, \emph{tini bhôy ônubhôb kôrlen}.

\bn{অনুভব} alone is the noun: what you feel from the inside.

\begin{cousinweb}
\textbf{\bn{অনুভব} \bn{করা}} (\emph{ônubhôb kôrā}) is Sanskrit \emph{anubhava}, \textquotedblleft{}perception, understanding\textquotedblright{} (Monier-Williams), with \emph{kôrā}, \textquotedblleft{}to do\textquotedblright{} (CDIAL 2814).
\end{cousinweb}

\begin{grammarlens}[title={What you've built}]
One more word for describing what you lived and felt.
\end{grammarlens}

\subsection*{Your turn}
\begin{itemize}
\raggedright
  \item \emph{Say it:} \emph{ônubhôb kôrā}
  \item \emph{Say it:} \emph{ônubhôb kôrā}, once more
  \item \emph{Say it:} \emph{mīrā bôndhur bæthā ônubhôb kôrlô.} — Mira felt her friend's pain.
  \item \emph{Recall:} say the Bengali for a feeling, then the Bengali for emotion, then say \emph{ônubhôb kôrā} again
\end{itemize}

\begin{tcolorbox}[breakable,colback=teal!4,colframe=teal!35!black,title={Before you move on}]
What does \bn{অনুভব} \bn{করা} mean? (\textquotedblleft{}To feel; to sense\textquotedblright{}.) Say it once more.
\end{tcolorbox}
\section[\texorpdfstring{\bn{টের} \bn{পাওয়া} (ṭer pāoyā)}{টের পাওয়া (ṭer pāoyā)}]{\bn{টের} \bn{পাওয়া} (ṭer pāoyā) — to become aware of; to sense}

\label{lesson:BN-C438-ter-paoya}



\begin{quote}
Before the new one: say the Bengali for emotion, then the Bengali for to feel.
\end{quote}

\subsection*{You'll want to know: \bn{টের} \bn{পাওয়া}}
\textbf{\bn{টের} \bn{পাওয়া}} — \emph{ṭer pāoyā} — \textquotedblleft{}to become aware of; to sense\textquotedblright{}.

\begin{quote}
\bn{রাতে} \bn{আমি} \bn{বাইরে} \bn{একটা} \bn{আওয়াজ} \bn{টের} \bn{পেলাম।}
\end{quote}

\begin{quote}
\emph{rāte āmi bāire ekṭā āoyāj ṭer pelām.}
\end{quote}

\begin{quote}
In the night I became aware of a noise outside.
\end{quote}

Past (he/she …): \bn{তিনি} \bn{বিপদ} \bn{টের} \bn{পেলেন}, \emph{tini bipôd ṭer pelen}.

\bn{টের} lives only inside \bn{টের} \bn{পাওয়া}, often for something you notice without seeing it.

\begin{cousinweb}
\textbf{\bn{টের} \bn{পাওয়া}} (\emph{ṭer pāoyā}) is \emph{ṭer}, \textquotedblleft{}notice\textquotedblright{}, which CDIAL lists for Bengali under a reconstructed form (CDIAL 5473), with \emph{pāoyā}, \textquotedblleft{}to get\textquotedblright{} (CDIAL 8943).
\end{cousinweb}

\begin{grammarlens}[title={What you've built}]
One more word for describing what you lived and felt.
\end{grammarlens}

\subsection*{Your turn}
\begin{itemize}
\raggedright
  \item \emph{Say it:} \emph{ṭer pāoyā}
  \item \emph{Say it:} \emph{ṭer pāoyā}, once more
  \item \emph{Say it:} \emph{rāte āmi bāire ekṭā āoyāj ṭer pelām.} — In the night I became aware of a noise outside.
  \item \emph{Recall:} say the Bengali for emotion, then the Bengali for to feel, then say \emph{ṭer pāoyā} again
\end{itemize}

\begin{tcolorbox}[breakable,colback=teal!4,colframe=teal!35!black,title={Before you move on}]
What does \bn{টের} \bn{পাওয়া} mean? (\textquotedblleft{}To become aware of; to sense\textquotedblright{}.) Say it once more.
\end{tcolorbox}
